\documentclass[11pt]{article}

\usepackage[margin=1in]{geometry}
\usepackage{amsmath,amssymb,mathtools}
\usepackage{booktabs}
\usepackage{tabularx}
\usepackage{enumitem}
\usepackage{xcolor}
\usepackage{microtype}
\usepackage{url}
\usepackage[hidelinks]{hyperref}

\providecommand{\reportbuilddatetime}{unknown}
\InputIfFileExists{lps_geometry_only_plateau_kd_selector_build_info.tex}{}{}

\newcommand{\R}{\mathbb{R}}
\newcommand{\Gfull}{\mathcal G_{\mathrm{full}}}
\newcommand{\Gplateau}{\mathcal G_{\mathrm{plateau}}}
\newcommand{\argmin}{\operatorname*{arg\,min}}
\newcommand{\CV}{\operatorname{CV}}
\newcommand{\TVE}{\operatorname{TVE}}
\newcommand{\RMSE}{\operatorname{RMSE}}
\newcommand{\code}[1]{\texttt{#1}}

\title{\textbf{Geometry-Only Plateau Selection for LPS}\\[3pt]
\large A support-size and chart-dimension selector for local PCA charts}
\author{\small Source:\\[-1pt]
\begin{minipage}{0.92\textwidth}
\centering\footnotesize\ttfamily
\detokenize{~/current_projects/geosmooth/dev/methods/lps/status/}\\
\detokenize{lps_geometry_only_plateau_kd_selector.tex}
\end{minipage}}
\date{Build \reportbuilddatetime}

\begin{document}
\maketitle

\section{Purpose}

This note defines the geometry-only \code{plateau\_kd} selector for local
polynomial smoothing (LPS).  The selector chooses a support size \(k\) and a
local PCA chart dimension \(d\) before looking at the response \(y\).  Its
purpose is to test whether a stable geometric dimension plateau can provide a
useful deployable alternative to response-CV selection over the full Cartesian
grid of \((k,d)\) candidates.

The selector is deliberately conservative in its first implementation.  It is
not intended to replace \code{full\_kd}; it is a new comparator.  The important
distinction is that \code{full\_kd} asks which \((k,d)\) has the best inner-CV
prediction score, while \code{plateau\_kd} asks which \((k,d)\) is suggested by
the input geometry alone.

\section{The local chart objects}

Let
\[
  X=\{x_1,\ldots,x_n\}\subset\R^p
\]
be the observed data cloud.  For an anchor \(a\), the support of size \(k\) is
\[
  U_a(k)=\{x_a\}\cup
  \{\text{the }k-1\text{ nearest neighbors of }x_a\}.
\]
The centered support matrix is
\[
  Z_a(k)
  =
  \begin{bmatrix}
  (x_i-x_a)^\top
  \end{bmatrix}_{x_i\in U_a(k)}.
\]
Let
\[
  \sigma_{a1}(k)\ge \sigma_{a2}(k)\ge\cdots\ge 0
\]
be the singular values of \(Z_a(k)\).  The total-variance-explained dimension at
threshold \(\tau\), currently \(\tau=0.95\), is
\[
  d_a(k;\tau)
  =
  \min\left\{
  d:
  \frac{\sum_{j=1}^d \sigma_{aj}(k)^2}
       {\sum_j \sigma_{aj}(k)^2}
  \ge \tau
  \right\}.
\]
This dimension is clipped to the feasible design range.  If the local
polynomial degree is \(g\), the number of polynomial columns is
\[
  q(d,g)=\binom{d+g}{g}.
\]
With design margin \(m\), a candidate is feasible only when
\[
  q(d,g)+m\le k.
\]
The current implementation also bounds \(d\) by the user-supplied numeric
\code{chart.dim.grid}; for example, the CSD comparisons use \(d\in\{1,\ldots,8\}\).
This keeps \code{plateau\_kd} inside the same numeric universe as
\code{full\_kd}.

\section{The plateau rule}

Let
\[
  \mathcal K=\{k_1<k_2<\cdots<k_L\}
\]
be the support-size grid.  For each representative anchor \(a\), compute the
dimension path
\[
  d_a(k_1;\tau),d_a(k_2;\tau),\ldots,d_a(k_L;\tau).
\]
The baseline dimension is
\[
  d_a^0=d_a(k_1;\tau).
\]
The plateau endpoint is the last support size in the initial stable prefix:
\[
  k_a^\star
  =
  \max\left\{
  k_\ell:
  d_a(k_r;\tau)=d_a^0
  \text{ for every } r\le \ell
  \right\}.
\]
This is an initial-prefix rule.  If the dimension changes and later returns to
\(d_a^0\), the later return is ignored.  This makes the rule closer to the
simple geometric question: how far can the neighborhood grow before the local
dimension estimate first changes?

For a global LPS fit, the anchor-specific quantities are aggregated by medians:
\[
  k^\star
  =
  \operatorname{nearest}_{k\in\mathcal K}
  \operatorname{median}_a k_a^\star,
  \qquad
  d^\star
  =
  \operatorname{clip}
  \left(
  \operatorname{round}\operatorname{median}_a d_a^0
  \right).
\]
The clipping enforces the supplied dimension grid, the ambient dimension, and
the polynomial-design feasibility inequality at \(k^\star\).  The resulting
candidate family is therefore
\[
  \Gplateau=\{(k^\star,d^\star)\},
\]
with one row for each requested degree, kernel, and bandwidth multiplier.  In
the CSD comparison reported with this note, the degree, kernel, and bandwidth
multiplier are fixed, so \(\Gplateau\) contains one candidate.

\section{Comparison with existing selectors}

\begin{tabularx}{\textwidth}{@{}lX@{}}
\toprule
Selector & Meaning \\
\midrule
\code{auto} &
Uses observed \(X\) to estimate one global chart dimension \(d\), then selects
support size \(k\) by the ordinary LPS CV grid.  The chosen dimension is shared
by all anchors. \\
\code{local.auto} &
Uses observed \(X\) to estimate a dimension \(d_a\) separately for each anchor.
The support size is still selected by the ordinary LPS CV grid. \\
\code{full\_kd} &
Evaluates the full Cartesian numeric grid
\(\Gfull=\mathcal K\times\mathcal D\) by inner CV and chooses the one global
pair \((k^\star,d^\star)\) with the best CV score. \\
\code{sparse\_kd} &
Evaluates a smaller deterministic skeleton of the same Cartesian grid by inner
CV.  It is a response-CV shortcut to \code{full\_kd}. \\
\code{plateau\_kd} &
Uses \(X\) only to choose one global \((k^\star,d^\star)\) from the initial
dimension plateau.  CV is still reported for the selected candidate, but CV
does not choose the pair. \\
\bottomrule
\end{tabularx}

\section{What the comparison should reveal}

The \code{plateau\_kd} selector tests a specific hypothesis.  If the dominant
mistake in \code{full\_kd} is noisy inner-CV ranking over many nearby
candidates, then a geometry-only plateau may sometimes perform comparably while
being much cheaper and more stable.  If the response really needs a support
scale different from the geometry plateau, then \code{plateau\_kd} should lose
to \code{full\_kd}, and the loss should be visible in truth-facing RMSE.

The associated HTML comparison therefore reports, for each synthetic dataset,
\[
  R_m
  =
  \left\{
  \frac{1}{n}\sum_{i=1}^n
  \left(\widehat f_m(x_i)-f_i\right)^2
  \right\}^{1/2},
\]
where \(f_i\) is the known synthetic truth and \(\widehat f_m\) is the fitted
LPS estimate under method \(m\).  The report also records the selected support
size, selected chart dimension, observed CV score, and runtime.  In the 1D and
2D examples the truth and fitted functions are plotted directly; in higher
dimensions the report uses summary tables because direct surface display is not
meaningful.

\section{Limitations}

The current \code{plateau\_kd} implementation is intentionally a first,
auditable rule.  It uses the total-variance-explained dimension only; it does
not yet combine variance thresholds with eigengaps, participation rank, or a
response-facing near-tie rule.  It aggregates local plateau endpoints into one
global \((k,d)\), so it is not an anchor-specific local-scale method.  A later
local variant could use \(k_a^\star\) and \(d_a^0\) anchor by anchor, but that
would be a different selector and should be tested separately.

\section{Implementation status}

The selector is implemented through the shared coupled candidate-specification
path.  The public trigger is
\[
  \code{selection.strategy = "plateau\_kd"}
\]
with \(\code{coordinate.method = "local.pca"}\) and an explicit numeric
\code{chart.dim.grid}.  The fit object records the strategy, the selected
candidate, and per-anchor plateau diagnostics under
\[
  \code{fit\$diagnostics\$coupled.kd.selection}.
\]
The implementation includes a unit test that verifies the one-candidate LPS path
and verifies that the selector refuses to run without an explicit numeric
dimension universe.

\end{document}
